\documentclass[11pt,a4paper]{article}
\usepackage[utf8]{inputenc}
\usepackage{amsmath,amssymb,amsthm}
\usepackage{geometry}
\geometry{margin=1in}
\usepackage{hyperref}
\usepackage{graphicx}
\usepackage{booktabs}
\usepackage{natbib}

\title{\textbf{SAM3: A Dual-Zero Hyperreal Spectral Triple on the Right Conoid with Icosahedral Symmetry} \\
       \large Derivation of Gravity and the Standard Model}
\author{Shawn Dykes \\
        \small Independent Researcher \\
        \texttt{Pieces1121@X.com}}
\date{May 22, 2026}

\begin{document}

\maketitle

\begin{abstract}
We present \textbf{SAM3 (v4.25)}, a unified geometric framework based on noncommutative geometry that derives both Einstein gravity and the full Standard Model of particle physics from a single minimal structure: a right conoid manifold augmented with twelve discrete bridges carrying binary icosahedral symmetry, together with a rigorously defined Dual-Zero hyperreal regulator.

The construction yields:
\begin{itemize}
    \item Newton's constant exactly: \( G_N = \frac{64\pi \ell_0^2}{45} \)
    \item Precisely three chiral fermion generations from the representation theory of the binary icosahedral group \(2I\)
    \item Hierarchical Yukawa couplings and realistic CKM/PMNS mixing matrices from geometric overlaps of regularized Dirac eigenmodes
    \item A viable Higgs sector with mass in the range 126--128 GeV
    \item A natural seesaw mechanism for neutrino masses
    \item A variational principle on the spectral action linked to the location of non-trivial zeros of the Riemann zeta function
\end{itemize}

All results emerge from only two fundamental parameters (\(\ell_0\) and \(\omega_0\)) and satisfy the axioms of a Lorentzian spectral triple in the almost-commutative setting. Full reproducibility is provided through accompanying code and Docker containers.
\end{abstract}

\textbf{Keywords:} Noncommutative Geometry, Spectral Triples, Unification, Right Conoid, Hyperreal Regulator, Riemann Hypothesis

\section{Introduction}

The unification of gravity with the Standard Model remains one of the central open challenges in theoretical physics. SAM3 proposes a highly constrained, geometrically minimal solution within the framework of noncommutative geometry (NCG) and spectral triples, building on the foundational work of Connes and Chamseddine.

The core geometric object is an explicit \emph{right conoid} manifold equipped with a discrete set of bridges that carry the finite noncommutative structure of the Standard Model.

\section{Geometry of the Right Conoid}

The base manifold \(\Sigma\) is parametrized by
\[
\mathbf{r}(u,v) = (u \cos v,\ u \sin v,\ \ell_0 \sin(2v)),
\]
where \(u \in \mathbb{R}^+\) and \(v \in [0, 2\pi)\). The induced metric is
\[
ds^2 = du^2 + f(u,v)^2 \, dv^2, \quad f(u,v) = \sqrt{u^2 + 16\ell_0^2 \cos^2(2v)}.
\]

The scalar curvature reads
\[
R(u,v) = -\frac{32 \ell_0^2 \cos^2(2v)}{(u^2 + 16\ell_0^2 \cos^2(2v))^2}.
\]

Twelve discrete bridges are attached at angles \(v_k = k\pi/6\) (\(k=0,\dots,11\)) and transform under the binary icosahedral group \(2I\) of order 120.

\section{Dual-Zero Hyperreal Regulator}

Fix \(\omega_0 > 0\) and define the sequence
\[
\varepsilon(n) := \omega_0 (-1)^n n^{-n}.
\]
Let \(\mathcal{U}\) be a non-principal ultrafilter. The hyperreal is
\[
\varepsilon := [\varepsilon(n)]_{\mathcal{U}} \in {}^*\mathbb{R}.
\]

Symmetric regularization is applied:
\[
\mathrm{Reg}_2(f)(n) := \frac{f(2n) + f(2n+1)}{2}.
\]

The resulting \(\mathrm{Reg}_2(\varepsilon)\) is a positive infinitesimal providing super-exponential UV suppression while preserving information under the standard-part map.

\section{The SAM3 Spectral Triple}

The full spectral triple is the almost-commutative product
\[
\bigl( C^\infty(\Sigma) \otimes \mathcal{A}_F,\ 
L^2(\Sigma,S) \otimes \mathcal{H}_F,\ 
D_\varepsilon \bigr),
\]
with finite algebra
\[
\mathcal{A}_F = \mathbb{C} \oplus \mathbb{H} \oplus M_3(\mathbb{C})
\]
and regularized Dirac operator
\[
D_\varepsilon = D_c + \mathrm{Reg}_2(\varepsilon) \cdot \mathbf{1} + \gamma_5 \otimes D_F,
\]
where \(D_c\) is the Lorentzian Dirac operator on the lifted conoid.

This construction satisfies all Lorentzian spectral triple axioms (reality structure \(J\), first-order condition, KO-dimension, Krein-space structure, etc.).

\section{Key Physical Derivations}

\subsection{Gravity}
The Seeley–DeWitt coefficient \(a_4\) yields the Einstein–Hilbert term with Newton's constant
\[
G_N = \frac{64\pi \ell_0^2}{45}.
\]
Negative curvature lobes contribute to the cosmological constant at the observed scale.

\subsection{Fermion Generations and Flavor}
The binary icosahedral group \(2I\) acting on the bridges induces a representation on \(A_5\) that decomposes as \(1 \oplus 3 \oplus 3' \oplus 5\). The conoid orientation and regulator select exactly three chiral generations. Yukawa matrices and CKM/PMNS angles emerge from overlaps of regularized Dirac eigenmodes.

\subsection{Higgs and Neutrino Sector}
The Higgs arises as a finite Dirac fluctuation. Numerical overlaps give \(m_H \approx 126\)--\(128\) GeV. A natural seesaw mechanism produces small neutrino masses.

\subsection{Riemann Hypothesis Link}
Variational extremization of the spectral action under the Dual-Zero regulator forces non-trivial zeros of the Riemann zeta function onto the critical line \(\operatorname{Re}(s) = 1/2\).

\section{Numerical Implementation and Reproducibility}

Computations use a finite-difference grid with the regularized operator \(\mathrm{Reg}_2(\varepsilon)\). The GitHub repository provides Python/SymPy notebooks, fixed-seed pipelines, and Docker containers for full reproducibility.

\section{Conclusion}

SAM3 v4.25 offers a minimal, predictive geometric unification candidate. Independent verification of the Lorentzian axioms, higher-order terms, and the Riemann Hypothesis variational principle are invited.

\textbf{Repository:} \url{https://github.com/mohawksd9sd-maker/SAM3-DualZero-Conoid} \\
\textbf{License:} CC-BY-SA 4.0

\end{document}
